\section{Conclusions}

We introduced and studied the \textsc{Generalized Conductance}, a problem important from the perspective of potential applications. Our main algorithm combines two complementary reductions which yield an $\bigo\br{\log n}$-approximation on general graphs and an $\bigo\br{1}$-approximation on trees. For multiplicative demand functions, we further improve the approximation guarantee to $\bigo\br{\sqrt{\log n}}$.

Beyond the core objective, we showed that these approximation ratios carry over to natural downstream tasks. In particular, we obtained an $\bigo\br{\log n}$ bicriteria approximation for \textsc{Graph Partitioning with Demands} and an $\bigo\br{\log n}$-approximation for \textsc{Hierarchical Clustering with Demands}, with corresponding $\bigo\br{\sqrt{\log n}}$ guarantees in the multiplicative setting and $\bigo\br{1}$-approximation for trees. These consequences indicate that generalized conductance can serve as a robust algorithmic primitive, similarly to how sparsest-cut-type objectives are used in classical divide-and-conquer methods. It is of interest whether our algorithm is applicable to different tasks in the presence of demands, such as network design or routing.

It is also interesting whether our generalized setup can also be incorporated into the framework of spectral graph theory. In particular, it is known that the classical conductance is closely related to the second eigenvalue of the Laplacian matrix of a graph. It would be interesting to see whether a similar relationship can be established for generalized conductance and some appropriate generalization of the Laplacian matrix.
